\chapter{Практика}

\section{Построение модели}

Построим модель несколько раз с разными значениями max\_depth и max\_features и вычислим MSE оценку ошибки, чтобы найти оптимальные значения для модели:

\begin{figure}[H]
    \centering
    \includegraphics{10.png}
    \caption{MSE для max\_depth}
\end{figure}

\begin{figure}[H]
    \centering
    \includegraphics{20.png}
    \caption{MSE для max\_features}
\end{figure}

Отчет по модели

\begin{figure}[H]
    \centering
    \includegraphics{40.png}
    \caption{Отчет по модели}
\end{figure}
 
Оптимальные значения max\_depth=4, max\_features = 15. \\

Построим итоговую модель с параметрами выше:

\begin{figure}[H]
    \centering
    \includegraphics{30.png}
    \caption{Дерево}
\end{figure}

По дереву можно определить логику, по которой модель будет принимать решения.

\endinput